\clearpage

\section{OBJETIVOS DEL PROYECTO}

El presente capítulo establece los objetivos que orientan el desarrollo del proyecto, derivados directamente del problema identificado y de los componentes funcionales que estructuran la propuesta. Se define un objetivo general que expresa el propósito central del trabajo, acompañado de objetivos específicos que lo descomponen en elementos técnicos concretos y abordables de manera independiente.

\vspace{0.5em}
\subsection{Objetivo general}

\vspace{0.3em}
Implementar una plataforma web para la administración del ciclo completo del Desafío Bebras en Bolivia, orientada a la gestión de unidades educativas escolares participantes, el registro de estudiantes de nivel primario y secundario, y la 
ejecución organizada de competencias de pensamiento computacional, que permitan su realización estructurada y continua en el contexto nacional.

\vspace{0.5em}
\subsection{Objetivos específicos}

\begin{itemize}
    \item Implementar una arquitectura modular de la plataforma web, definiendo la estructura tecnológica que integra y sostiene los distintos componentes del sistema de manera cohesionada y extensible.

    \item Implementar un módulo de gestión de competencias que permita la creación, categorización y administración de tareas de pensamiento computacional según rangos de edad y niveles de dificultad establecidos por la organización internacional.

    \item Implementar un sistema de gestión de contenido (CMS) que permita la administración dinámica de información institucional, publicaciones y estructura del sitio web informativo del Desafío Bebras Bolivia, orientado a docentes e 
    instituciones escolares.

     \item Implementar un módulo de registro y administración de participantes que permita la inscripción de unidades educativas escolares, docentes responsables y estudiantes de nivel primario y secundario en el concurso.

    \item Implementar mecanismos de evaluación que soporten la resolución de tareas en entorno digital y la generación de formatos imprimibles para su procesamiento mediante reconocimiento óptico de caracteres, adaptados a las condiciones de conectividad de unidades educativas escolares bolivianas.
\end{itemize}